\section{Cohomology and base change, I}
\label{section-cohomology-and-base-change}

\noindent
Let $f : X \to S$ be a morphism of schemes.
Let $\mathcal{F}$ be a quasi-coherent sheaf on $X$.
Suppose further that $g : S' \to S$ is any morphism of schemes. Denote
$X' = X_{S'} = S' \times_S X$ the base change of $X$ and denote
$f' : X' \to S'$ the base change of $f$.
Also write $g' : X' \to X$ the projection,
and set $\mathcal{F}' = (g')^*\mathcal{F}$.
Here is a diagram representing the situation:
\begin{equation}
\label{equation-base-change-diagram}
\vcenter{
\xymatrix{
\mathcal{F}' = (g')^*\mathcal{F} &
X' \ar[r]_{g'} \ar[d]_{f'} &
X \ar[d]^f &
\mathcal{F} \\
Rf'_*\mathcal{F}' &
S' \ar[r]^g &
S &
Rf_*\mathcal{F}
}
}
\end{equation}
Here is the simplest case of the base change property we have in mind.

\begin{lemma}
\label{lemma-affine-base-change}
Let $f : X \to S$ be a morphism of schemes.
Let $\mathcal{F}$ be a quasi-coherent $\mathcal{O}_X$-module.
Assume $f$ is affine.
In this case $f_*\mathcal{F} \cong Rf_*\mathcal{F}$ is
a quasi-coherent sheaf, and for every base change diagram
(\ref{equation-base-change-diagram})
we have
$$
g^*f_*\mathcal{F} = f'_*(g')^*\mathcal{F}.
$$
\end{lemma}

\begin{proof}
The vanishing of higher direct images is
Lemma \ref{lemma-relative-affine-vanishing}.
The statement is local on $S$ and $S'$. Hence we may
assume $X = \Spec(A)$, $S = \Spec(R)$,
$S' = \Spec(R')$ and $\mathcal{F} = \widetilde{M}$
for some $A$-module $M$.
We use Schemes, Lemma \ref{schemes-lemma-widetilde-pullback}
to describe pullbacks and pushforwards of $\mathcal{F}$.
Namely, $X' = \Spec(R' \otimes_R A)$ and
$\mathcal{F}'$ is the quasi-coherent sheaf associated
to $(R' \otimes_R A) \otimes_A M$.
Thus we see that the lemma boils down to the
equality
$$
(R' \otimes_R A) \otimes_A M = R' \otimes_R M
$$
as $R'$-modules.
\end{proof}

\noindent
In many situations it is sufficient to know about the following
special case of cohomology and base change. It follows immediately
from the stronger results in
Section \ref{section-cohomology-and-base-change-derived},
but since it is so important it deserves its own proof.

\begin{lemma}[Flat base change]
\label{lemma-flat-base-change-cohomology}
Consider a cartesian diagram of schemes
$$
\xymatrix{
X' \ar[d]_{f'} \ar[r]_{g'} & X \ar[d]^f \\
S' \ar[r]^g & S
}
$$
Let $\mathcal{F}$ be a quasi-coherent $\mathcal{O}_X$-module
with pullback $\mathcal{F}' = (g')^*\mathcal{F}$.
Assume that $g$ is flat and that $f$ is quasi-compact and quasi-separated.
For any $i \geq 0$
\begin{enumerate}
\item the base change map of
Cohomology, Lemma \ref{cohomology-lemma-base-change-map-flat-case}
is an isomorphism
$$
g^*R^if_*\mathcal{F} \longrightarrow R^if'_*\mathcal{F}',
$$
\item if $S = \Spec(A)$ and $S' = \Spec(B)$, then
$H^i(X, \mathcal{F}) \otimes_A B = H^i(X', \mathcal{F}')$.
\end{enumerate}
\end{lemma}

\begin{proof}
Using Cohomology, Lemma \ref{cohomology-lemma-base-change-map-flat-case} in (1)
is allowed since $g'$ is flat by
Morphisms, Lemma \ref{morphisms-lemma-base-change-flat}.
Having said this, part (1) follows from part (2). Namely,
part (1) is local on $S'$ and hence we may assume $S$
and $S'$ are affine. In other words, we have $S = \Spec(A)$
and $S' = \Spec(B)$ as in (2).
Then since $R^if_*\mathcal{F}$ is quasi-coherent
(Lemma \ref{lemma-quasi-coherence-higher-direct-images}),
it is the quasi-coherent $\mathcal{O}_S$-module associated to the
$A$-module $H^0(S, R^if_*\mathcal{F}) = H^i(X, \mathcal{F})$
(equality by
Lemma \ref{lemma-quasi-coherence-higher-direct-images-application}).
Similarly, $R^if'_*\mathcal{F}'$ is the quasi-coherent
$\mathcal{O}_{S'}$-module associated to the $B$-module
$H^i(X', \mathcal{F}')$. Since pullback by $g$ corresponds
to $- \otimes_A B$ on modules
(Schemes, Lemma \ref{schemes-lemma-widetilde-pullback})
we see that it suffices to prove (2).

\medskip\noindent
Let $A \to B$ be a flat ring homomorphism.
Let $X$ be a quasi-compact and quasi-separated scheme over $A$.
Let $\mathcal{F}$ be a quasi-coherent $\mathcal{O}_X$-module.
Set $X_B = X \times_{\Spec(A)} \Spec(B)$ and denote
$\mathcal{F}_B$ the pullback of $\mathcal{F}$.
We are trying to show that the map
$$
H^i(X, \mathcal{F}) \otimes_A B \longrightarrow H^i(X_B, \mathcal{F}_B)
$$
(given by the reference in the statement of the lemma)
is an isomorphism.

\medskip\noindent
In case $X$ is separated, choose an affine open covering
$\mathcal{U} : X = U_1 \cup \ldots \cup U_t$ and recall that
$$
\check{H}^p(\mathcal{U}, \mathcal{F}) = H^p(X, \mathcal{F}),
$$
see
Lemma \ref{lemma-cech-cohomology-quasi-coherent}.
If $\mathcal{U}_B : X_B = (U_1)_B \cup \ldots \cup (U_t)_B$ we obtain
by base change, then it is still the case that each $(U_i)_B$ is affine
and that $X_B$ is separated. Thus we obtain
$$
\check{H}^p(\mathcal{U}_B, \mathcal{F}_B) = H^p(X_B, \mathcal{F}_B).
$$
We have the following relation between the {\v C}ech complexes
$$
\check{\mathcal{C}}^\bullet(\mathcal{U}_B, \mathcal{F}_B) =
\check{\mathcal{C}}^\bullet(\mathcal{U}, \mathcal{F}) \otimes_A B
$$
as follows from
Lemma \ref{lemma-affine-base-change}.
Since $A \to B$ is flat, the same thing remains true on taking cohomology.

\medskip\noindent
In case $X$ is quasi-separated, choose an affine open covering
$\mathcal{U} : X = U_1 \cup \ldots \cup U_t$. We will use the
{\v C}ech-to-cohomology spectral sequence
Cohomology, Lemma \ref{cohomology-lemma-cech-spectral-sequence}.
The reader who wishes to avoid this spectral sequence
can use Mayer-Vietoris and induction on $t$ as in the proof of
Lemma \ref{lemma-quasi-coherence-higher-direct-images}.
The spectral sequence has $E_2$-page
$E_2^{p, q} = \check{H}^p(\mathcal{U}, \underline{H}^q(\mathcal{F}))$
and converges to $H^{p + q}(X, \mathcal{F})$.
Similarly, we have a spectral sequence with $E_2$-page
$E_2^{p, q} = \check{H}^p(\mathcal{U}_B, \underline{H}^q(\mathcal{F}_B))$
which converges to $H^{p + q}(X_B, \mathcal{F}_B)$.
Since the intersections $U_{i_0 \ldots i_p}$ are quasi-compact
and separated, the result of the second paragraph of the proof gives
$\check{H}^p(\mathcal{U}_B, \underline{H}^q(\mathcal{F}_B)) =
\check{H}^p(\mathcal{U}, \underline{H}^q(\mathcal{F})) \otimes_A B$.
Using that $A \to B$ is flat we conclude that
$H^i(X, \mathcal{F}) \otimes_A B \to H^i(X_B, \mathcal{F}_B)$
is an isomorphism for all $i$ and we win.
\end{proof}

\begin{lemma}[Finite locally free base change]
\label{lemma-finite-locally-free-base-change-cohomology}
Consider a cartesian diagram of schemes
$$
\xymatrix{
Y \ar[d]_{g} \ar[r]_h & X \ar[d]^f \\
\Spec(B) \ar[r] & \Spec(A)
}
$$
Let $\mathcal{F}$ be a quasi-coherent $\mathcal{O}_X$-module
with pullback $\mathcal{G} = h^*\mathcal{F}$.
If $B$ is a finite locally free $A$-module, then
$H^i(X, \mathcal{F}) \otimes_A B = H^i(Y, \mathcal{G})$.
\end{lemma}

\noindent
{\bf Warning}: Do not use this lemma unless you understand the difference
between this and Lemma \ref{lemma-flat-base-change-cohomology}.

\begin{proof}
In case $X$ is separated, choose an affine open covering
$\mathcal{U} : X = \bigcup_{i \in I} U_i$ and recall that
$$
\check{H}^p(\mathcal{U}, \mathcal{F}) = H^p(X, \mathcal{F}),
$$
see
Lemma \ref{lemma-cech-cohomology-quasi-coherent}.
Let $\mathcal{V} : Y = \bigcup_{i \in I} g^{-1}(U_i)$
be the corresponding affine open covering of $Y$.
The opens $V_i = g^{-1}(U_i) = U_i \times_{\Spec(A)} \Spec(B)$
are affine and $Y$ is separated. Thus we obtain
$$
\check{H}^p(\mathcal{V}, \mathcal{G}) = H^p(Y, \mathcal{G}).
$$
We claim the map of {\v C}ech complexes
$$
\check{\mathcal{C}}^\bullet(\mathcal{U}, \mathcal{F}) \otimes_A B
\longrightarrow
\check{\mathcal{C}}^\bullet(\mathcal{V}, \mathcal{G})
$$
is an isomorphism. Namely, as $B$ is finitely presented as an $A$-module
we see that tensoring with $B$ over $A$ commutes with products, see
Algebra, Proposition \ref{algebra-proposition-fp-tensor}.
Thus it suffices to show that the maps
$\Gamma(U_{i_0 \ldots i_p}, \mathcal{F}) \otimes_A B \to
\Gamma(V_{i_0 \ldots i_p}, \mathcal{G})$ 
are isomorphisms which follows from
Lemma \ref{lemma-affine-base-change}.
Since $A \to B$ is flat, the same thing remains true on taking cohomology.

\medskip\noindent
In the general case we argue in exactly the same way using affine
open covering $\mathcal{U} : X = \bigcup_{i \in I} U_i$ and the
corresponding covering $\mathcal{V} : Y = \bigcup_{i \in I} V_i$
with $V_i = g^{-1}(U_i)$ as above. We will use the
{\v C}ech-to-cohomology spectral sequence
Cohomology, Lemma \ref{cohomology-lemma-cech-spectral-sequence}.
The spectral sequence has $E_2$-page
$E_2^{p, q} = \check{H}^p(\mathcal{U}, \underline{H}^q(\mathcal{F}))$
and converges to $H^{p + q}(X, \mathcal{F})$.
Similarly, we have a spectral sequence with $E_2$-page
$E_2^{p, q} = \check{H}^p(\mathcal{V}, \underline{H}^q(\mathcal{G}))$
which converges to $H^{p + q}(Y, \mathcal{G})$.
Since the intersections $U_{i_0 \ldots i_p}$ are separated, the result
of the previous paragraph gives isomorphisms
$\Gamma(U_{i_0 \ldots i_p}, \underline{H}^q(\mathcal{F})) \otimes_A B
\to \Gamma(V_{i_0 \ldots i_p}, \underline{H}^q(\mathcal{G}))$.
Using that $- \otimes_A B$ commutes with products and is exact, we conclude
that
$\check{H}^p(\mathcal{U}, \underline{H}^q(\mathcal{F})) \otimes_A B
\to \check{H}^p(\mathcal{V}, \underline{H}^q(\mathcal{G}))$
is an isomorphism. Using that $A \to B$ is flat we conclude that
$H^i(X, \mathcal{F}) \otimes_A B \to H^i(Y, \mathcal{G})$
is an isomorphism for all $i$ and we win.
\end{proof}









